% Generated draft: requires mathematical and source review.
\documentclass[11pt]{article}
\usepackage[margin=25mm]{geometry}
\usepackage{fontspec}
\setmainfont{Latin Modern Roman}
\usepackage{amsmath,amssymb,mathtools,unicode-math}
\setmathfont{Latin Modern Math}
\usepackage{enumitem,needspace,xurl,hyperref}
\setlength{\parindent}{0pt}
\setlength{\parskip}{5pt}
\setlength{\emergencystretch}{3em}
\newcommand{\N}{\mathbb{N}}
\newcommand{\Z}{\mathbb{Z}}
\newcommand{\Q}{\mathbb{Q}}
\newcommand{\R}{\mathbb{R}}
\newcommand{\C}{\mathbb{C}}
\newcommand{\F}{\mathbb{F}}
\newcommand{\A}{\mathbb{A}}
\newcommand{\PP}{\mathbb{P}}
\newcommand{\E}{\mathbb{E}}
\newcommand{\eps}{\varepsilon}
\newcommand{\dd}{\,\mathrm{d}}
\begin{document}
\section*{UnsolvedMath ID: 7000006}
\textbf{Source problem code:} catalogue entry
\subsection{Definitions and mathematical statement}
For a regular \(C^3\) closed space curve \(\gamma:S^1\to\R^3\), wherever \(\gamma'\times\gamma''\ne0\), its torsion is
\[
 \tau_\gamma(t)=\frac{\det(\gamma'(t),\gamma''(t),\gamma^{(3)}(t))}{|\gamma'(t)\times\gamma''(t)|^2}.
\]
Thus torsion is classically defined everywhere if the curve has nonvanishing curvature. Interpreting positive curvature of a surface in the usual Gaussian sense, the question asks: if a closed curve bounds a surface \(\Sigma\) in \(\R^3\) with \(K_\Sigma>0\), must there be at least four distinct points where \(\tau_\gamma=0\)? In a classical nonvanishing-curvature formulation the proposed conclusion is
\[
 \#\{t\in S^1:\tau_\gamma(t)=0\}\geq4.
\]
The curve regularity, handling of curvature-zero points, immersed versus embedded meaning of the curve and surface, and topology of \(\Sigma\) are not supplied in the short statement.

The registry reports a resolution by Ghomi for closed curves bounding \emph{simply connected locally convex} surfaces, including the corresponding positively curved case. The simply connected condition is explicit in that reported theorem but absent from the extracted question. Hence the registry's solved label should not be read here as a verified result for arbitrary topological types of bounding surfaces. The cited abstract attributes the question to Rosenberg, related to Yau's problem on boundaries of positively curved disks, and compares the result to Sedykh's four-vertex theorem. These are supplied attribution and literature reports, not independently established claims in this draft.
\subsection{Short English statement}
Must a closed space curve bounding a positively curved surface in \(\R^3\) have at least four torsion zeros? The supplied solution reference concerns simply connected locally convex bounding surfaces, a topology condition omitted from the short question.
\subsection{Sources}
\begin{itemize}
\item Local UnsolvedMath registry, permanent ID 7000006, source code catalogue entry; attributed to Mohammad Ghomi, proposed year 2019; update 2026-08-06; dated status solved.
\item Ghomi, \emph{Open Problems in Geometry of Curves and Surfaces} (2019), Problem 1.6, PDF page 7: \url{https://people.math.gatech.edu/~ghomi/Papers/op.pdf}. Registry access date 2026-07-29; normalized PDF extraction; public author-hosted PDF; current status and statement reuse terms require review.
\item M. Ghomi, ``Boundary torsion and convex caps of locally convex surfaces'', arXiv:1501.07626; reported in the registry as published in \emph{Journal of Differential Geometry} (2017). The quoted abstract states the four-torsion-zero result for simply connected locally convex bounding surfaces and credits Rosenberg's question, Yau's related problem, and Sedykh's earlier theorem. Exact hypotheses and bibliographic details have not been independently checked.
\end{itemize}
\end{document}
